\chapter{Discussion}
\section{L'échantillonage permet-il de répondre à la problématique ?}
- biais d'échantillonage et régions/années moins représentées => perf du modèle diminuées dans certaine régions
- distribution shift => un modèle n'est pas sufisant pour prédire toute la ROI (sauf peut être s'il apprend la causalité).
- valeurs manquantes MNAR et biais de sélection (couverture nuageuse)

\section{Les variables mesurées représentent elles correctement les construits ?}
incertitudes liées aux données :
- NASC et acoustique active
- classes de phyto a partir des pigments

\section{Prédiction ou explicabilité ?}

\section{Contrainte des modèles prédictifs}
- nécessité d'avoir une métrique rigoureuse d'évaluation des performances du modèle (CV blocage)
- nécessité de robustesse façe à l'extrapolation 
- en découle la nécessité de limiter le shortcut learning
- apprendre la causalité permet la robustesse face à extrapolation 

\section{Pertinence des forêts aléatoires}
- surapprentissage et RF : surapprentissage possible par profondeur des arbres
- diminution de la variance utile pour le trade-off biais-variance de la prédiction
- extrapolation difficile étant donné que RF est fonction continue par morceau 
- importance des variables sur echantillon OOB mais OOB est leaked/fortement correlé/biais d'echantillonage

\section{Mise en relation des résultats associatifs avec la théorie de causalité et hypothèses}
- liens entre mésoéchelle et distrib phyto 
- liens trophiques entre phyto et NTI mais délai
- liens avec importance des variables
- liens entre connaissance des distribs de euphausiaces et prédictions
- liens entre connaissance inversion 18/38 et observations + distributions des poissons 38 et prédictions

\section{Domaine d'applicabilité du modèles et précautions d'interprétations}

\section{perspectives}